%=============================
%  QUINTO CAPITULO
%=============================
\chapter{ESTABILIDAD, CAOS Y LA JERARQUÍA DE ANILLOS DE FOTONES}

\section{El factor de adelgazamiento de Sneppen \texorpdfstring{$f_{thin} = e^{-2\pi\gamma}$}{f\_thin = e^{-2 pi gamma}}}
\citet{Sneppen} demostró que los fotones atrapados momentáneamente cerca de la esfera de fotones forman una jerarquía fractal de sub-anillos \citep{Sneppen, pertbuciiones_teoria_en_bh}. Cada sub-anillo $n$ corresponde a fotones que completan $n$ medias órbitas. La relación de espesores entre sub-anillos consecutivos se rige por el factor de adelgazamiento de Sneppen \citep{Sneppen}:
\begin{equation}
    f_{thin} = e^{-2\pi\gamma_\phi}
    \label{eq:factor_sneppen}
\end{equation}
En el caso estático de Schwarzschild, $\gamma_\phi = 1$, lo que da $f_{thin} = e^{-2\pi} \approx 1/535.5$ \citep{Sneppen, chaos_fractals_photon_ring}.

\section{Ensanchamiento de los anillos por rotación y carga}
En el espaciotiempo de Kerr-Newman, la rotación degrada el exponente de Lyapunov para órbitas progradas ($\gamma_\phi < 1$), reduciendo el factor de adelgazamiento de 535 a valores tan bajos como $f_{thin} \approx 2$ \citep{Sneppen, bh_mergers_photon_ring}. Esto ensancha los sub-anillos $n \ge 1$, permitiendo que sean detectados por instrumentos como el ngEHT y BHEX \citep{Measuring_Photon_Rings_with_the_ngEHT, 2406.09498v1}.

\section{Transición al caos en el espacio de fases}
Mientras que Kerr-Newman es integrable debido a la constante de Carter \citep{1403.5234v2, Achour_2025_J._Cosmol._Astropart._Phys._2025_012}, la presencia de perturbaciones no integrables destruye los toros de Kolmogorov-Arnold-Moser (KAM), generando caos determinista en el espacio de fases \citep{chaos_fractals_photon_ring}.
\label{cap5}
